\renewcommand{\tema}{Examen Diagnóstico - Cinemática}

\twocolumn
\section{Instrucciones}

Dispones de \textbf{15 minutos} para resolver 10 preguntas.

\begin{instruccionbox}
\subsection{Sobre la velocidad de respuesta:}

El examen de ingreso a la UNAM te da aproximadamente 1.5 minutos por pregunta. La experiencia de miles de aspirantes muestra que reprobar no siempre es por falta de conocimiento, sino por no terminar el examen. Saber contestar rápido es tan importante como conocer el tema.

Este diagnóstico replica las condiciones de exámenes estandarizados de certificación profesional. Evaluamos tres aspectos:

\begin{itemize}
    \item Conocimiento del tema
    \item Velocidad de respuesta
    \item Capacidad para eliminar distractores
\end{itemize}
\end{instruccionbox}

\subsection{Procedimiento:}

\begin{enumerate}
    \item Responde las 10 preguntas en máximo 15 minutos
    \item Si una pregunta te toma más de 2 minutos, marca cualquier respuesta y continúa
    \item Calcula mentalmente, sin calculadora
    \item Anota tus respuestas en formato: 1-A, 2-C, 3-B...
    \item Al terminar, revisa las respuestas en la siguiente sección
\end{enumerate}

\subsection{Propósito del diagnóstico:}

Este examen mide tu punto de partida. No importa cuántas respondas correctamente. El objetivo es identificar qué necesitas reforzar.

\vspace{1cm}

\begin{tcolorbox}[colback=amarillonotа, colframe=orange, boxrule=0pt, leftrule=3pt]
\textbf{Inicia el cronómetro y comienza.}
\end{tcolorbox}

\newpage

\section{PREGUNTAS DIAGNÓSTICO}

\begin{preguntabox}
\subsection{Pregunta 1}
Un automóvil recorre 120 km en 2 horas manteniendo su velocidad constante. ¿Cuál es su rapidez?

\begin{enumerate}
    \item[A)] 40 km/h
    \item[B)] 60 km/h
    \item[C)] 80 km/h
    \item[D)] 120 km/h
\end{enumerate}
\end{preguntabox}

\begin{preguntabox}
\subsection{Pregunta 2}
¿Cuál de las siguientes magnitudes es vectorial?

\begin{enumerate}
    \item[A)] Rapidez
    \item[B)] Distancia
    \item[C)] Tiempo
    \item[D)] Velocidad
\end{enumerate}
\end{preguntabox}

\begin{preguntabox}
\subsection{Pregunta 3}
Un objeto se mueve con velocidad constante. ¿Cuál afirmación describe correctamente su aceleración?

\begin{enumerate}
    \item[A)] Es positiva
    \item[B)] Es negativa
    \item[C)] Es cero
    \item[D)] Es variable
\end{enumerate}
\end{preguntabox}

\begin{preguntabox}
\subsection{Pregunta 4}
Un tren parte del reposo y alcanza 20 m/s en 10 segundos. ¿Cuál es su aceleración?

\begin{enumerate}
    \item[A)] 0.5 m/s²
    \item[B)] 2 m/s²
    \item[C)] 10 m/s²
    \item[D)] 200 m/s²
\end{enumerate}
\end{preguntabox}

\newpage

\begin{preguntabox}
\subsection{Pregunta 5}
En una gráfica posición vs tiempo, una línea recta horizontal indica que el objeto:

\begin{enumerate}
    \item[A)] Se mueve con velocidad constante
    \item[B)] Se mueve con aceleración constante
    \item[C)] Está en reposo
    \item[D)] Se mueve cada vez más rápido
\end{enumerate}
\end{preguntabox}

\begin{preguntabox}
\subsection{Pregunta 6}
Un objeto cae libremente desde el reposo. Después de 3 segundos, ¿qué velocidad tiene?\\
Considera $g = 10$ m/s²

\begin{enumerate}
    \item[A)] 3 m/s
    \item[B)] 10 m/s
    \item[C)] 30 m/s
    \item[D)] 90 m/s
\end{enumerate}
\end{preguntabox}

\begin{preguntabox}
\subsection{Pregunta 7}
La siguiente gráfica representa la velocidad de un objeto en función del tiempo. ¿Qué tipo de movimiento describe?

\begin{center}
\begin{tikzpicture}
\begin{axis}[
    width=10cm,
    height=7cm,
    axis lines=middle,
    xlabel={$t$},
    ylabel={$v$},
    xmin=0, xmax=5,
    ymin=0, ymax=5,
    xtick=\empty,
    ytick=\empty,
    thick,
    every axis x label/.style={at={(current axis.right of origin)},anchor=west},
    every axis y label/.style={at={(current axis.above origin)},anchor=south}
]
\addplot[color=azuloscuro, very thick, domain=0:4] {x};
\end{axis}
\end{tikzpicture}
\end{center}

\begin{enumerate}
    \item[A)] Movimiento rectilíneo uniforme
    \item[B)] Movimiento con velocidad constante
    \item[C)] Movimiento uniformemente acelerado
    \item[D)] Objeto en reposo
\end{enumerate}
\end{preguntabox}

\begin{preguntabox}
\subsection{Pregunta 8}
Se lanza un proyectil desde el origen con velocidad inicial de 20 m/s formando un ángulo de 30° con la horizontal. ¿Cuál es la altura máxima que alcanza?\\
Considera $g = 10$ m/s² y $\sin(30°) = 0.5$

\begin{enumerate}
    \item[A)] 5 m
    \item[B)] 10 m
    \item[C)] 20 m
    \item[D)] 40 m
\end{enumerate}
\end{preguntabox}

\begin{preguntabox}
\subsection{Pregunta 9}
En una gráfica velocidad vs tiempo para un objeto en caída libre, la pendiente de la recta representa:

\begin{enumerate}
    \item[A)] La posición del objeto
    \item[B)] La distancia recorrida
    \item[C)] La aceleración gravitacional
    \item[D)] La velocidad inicial
\end{enumerate}
\end{preguntabox}

\begin{preguntabox}
\subsection{Pregunta 10}
Un proyectil se lanza horizontalmente desde una altura de 20 m con velocidad de 10 m/s. ¿Cuánto tiempo tarda en tocar el suelo?\\
Considera $g = 10$ m/s²

\begin{enumerate}
    \item[A)] 1 s
    \item[B)] 2 s
    \item[C)] 4 s
    \item[D)] 10 s
\end{enumerate}
\end{preguntabox}

\vspace{1cm}

\begin{tcolorbox}[colback=amarillonotа, colframe=orange, boxrule=0pt, leftrule=3pt]
\textbf{Detén el cronómetro. Pasa a la siguiente sección para revisar tus respuestas.}
\end{tcolorbox}

\newpage
\onecolumn
\section{RESPUESTAS Y RETROALIMENTACIÓN}

\begin{respuestabox}
\subsection{Pregunta 1}
\textcolor{verdecorrecto}{\textbf{Respuesta correcta: B) 60 km/h}}

\textbf{Justificación:}\\
La rapidez se obtiene dividiendo la distancia entre el tiempo:
$$v = \frac{d}{t} = \frac{120 \text{ km}}{2 \text{ h}} = 60 \text{ km/h}$$

\textbf{Respuestas incorrectas:}
\begin{itemize}
    \item \textbf{A) 40 km/h:} Error en la operación aritmética. $\frac{120}{2} = 60$, no 40.
    \item \textbf{C) 80 km/h:} Este valor no resulta de ninguna operación válida con los datos proporcionados.
    \item \textbf{D) 120 km/h:} Confunde la distancia total con la rapidez. Se debe dividir entre el tiempo.
\end{itemize}
\end{respuestabox}

\begin{respuestabox}
\subsection{Pregunta 2}
\textcolor{verdecorrecto}{\textbf{Respuesta correcta: D) Velocidad}}

\textbf{Justificación:}\\
La velocidad requiere magnitud, dirección y sentido. Por ejemplo: ``80 km/h hacia el norte'' especifica completamente el vector $\vec{v}$.

\textbf{Respuestas incorrectas:}
\begin{itemize}
    \item \textbf{A) Rapidez:} Solo tiene magnitud. No indica dirección.
    \item \textbf{B) Distancia:} Magnitud escalar que indica cuánto se recorrió, sin especificar dirección.
    \item \textbf{C) Tiempo:} Magnitud escalar que mide duración. No tiene dirección asociada.
\end{itemize}
\end{respuestabox}

\begin{respuestabox}
\subsection{Pregunta 3}
\textcolor{verdecorrecto}{\textbf{Respuesta correcta: C) Es cero}}

\textbf{Justificación:}\\
La aceleración mide el cambio de velocidad: $a = \frac{\Delta v}{\Delta t}$. Si la velocidad es constante, $\Delta v = 0$, entonces $a = 0$.

\textbf{Respuestas incorrectas:}
\begin{itemize}
    \item \textbf{A) Es positiva:} Implicaría que la velocidad aumenta. Contradice la condición de velocidad constante.
    \item \textbf{B) Es negativa:} Significaría que el objeto está frenando. La velocidad constante excluye esta posibilidad.
    \item \textbf{D) Es variable:} Si la aceleración cambiara, la velocidad no sería constante.
\end{itemize}
\end{respuestabox}

\newpage

\begin{respuestabox}
\subsection{Pregunta 4}
\textcolor{verdecorrecto}{\textbf{Respuesta correcta: B) 2 m/s²}}

\textbf{Justificación:}\\
La aceleración es el cambio de velocidad dividido entre el tiempo:
$$a = \frac{v_f - v_i}{t} = \frac{20 - 0}{10} = 2 \text{ m/s}^2$$

\textbf{Respuestas incorrectas:}
\begin{itemize}
    \item \textbf{A) 0.5 m/s²:} Invierte la división. La operación correcta es $\frac{20}{10}$, no $\frac{10}{20}$.
    \item \textbf{C) 10 m/s²:} Usa solo el tiempo. Debe dividirse el cambio de velocidad entre el tiempo.
    \item \textbf{D) 200 m/s²:} Multiplica en lugar de dividir. La aceleración requiere $a = \frac{\Delta v}{\Delta t}$.
\end{itemize}
\end{respuestabox}

\begin{respuestabox}
\subsection{Pregunta 5}
\textcolor{verdecorrecto}{\textbf{Respuesta correcta: C) Está en reposo}}

\textbf{Justificación:}\\
Una línea horizontal indica que la posición no cambia con el tiempo. El objeto permanece en la misma ubicación: $x(t) = x_0 = \text{constante}$.

\textbf{Respuestas incorrectas:}
\begin{itemize}
    \item \textbf{A) Se mueve con velocidad constante:} El movimiento con velocidad constante genera una línea inclinada, no horizontal.
    \item \textbf{B) Se mueve con aceleración constante:} La aceleración constante produce una parábola en la gráfica posición-tiempo.
    \item \textbf{D) Se mueve cada vez más rápido:} La posición aumentaría progresivamente. La gráfica tendría pendiente creciente.
\end{itemize}
\end{respuestabox}

\begin{respuestabox}
\subsection{Pregunta 6}
\textcolor{verdecorrecto}{\textbf{Respuesta correcta: C) 30 m/s}}

\textbf{Justificación:}\\
En caída libre desde el reposo: $v = gt = (10 \text{ m/s}^2)(3 \text{ s}) = 30$ m/s. El objeto gana 10 m/s cada segundo.

\textbf{Respuestas incorrectas:}
\begin{itemize}
    \item \textbf{A) 3 m/s:} Usa solo el tiempo. La fórmula correcta es $v = gt$.
    \item \textbf{B) 10 m/s:} Usa solo la gravedad. Debe multiplicarse por el tiempo transcurrido.
    \item \textbf{D) 90 m/s:} Confunde la fórmula de velocidad ($v = gt$) con la de distancia ($d = \frac{1}{2}gt^2$).
\end{itemize}
\end{respuestabox}

\newpage

\begin{respuestabox}
\subsection{Pregunta 7}
\textcolor{verdecorrecto}{\textbf{Respuesta correcta: C) Movimiento uniformemente acelerado}}

\textbf{Justificación:}\\
Una línea recta con pendiente positiva en la gráfica velocidad-tiempo indica que la velocidad aumenta uniformemente. La pendiente representa la aceleración constante: $a = \frac{dv}{dt} = \text{constante}$.

\textbf{Respuestas incorrectas:}
\begin{itemize}
    \item \textbf{A) Movimiento rectilíneo uniforme:} El MRU requiere velocidad constante. En la gráfica, la velocidad aumenta.
    \item \textbf{B) Movimiento con velocidad constante:} La velocidad constante se representa con una línea horizontal en esta gráfica, no inclinada.
    \item \textbf{D) Objeto en reposo:} Un objeto en reposo tendría velocidad cero en todo momento. La gráfica mostraría una línea sobre el eje del tiempo.
\end{itemize}
\end{respuestabox}

\begin{respuestabox}
\subsection{Pregunta 8}
\textcolor{verdecorrecto}{\textbf{Respuesta correcta: A) 5 m}}

\textbf{Justificación:}\\
La altura máxima se calcula con:
$$h_{\text{max}} = \frac{(v_0 \sin \theta)^2}{2g}$$

Velocidad vertical inicial: $v_{0y} = v_0 \sin \theta = 20 \times 0.5 = 10$ m/s

$$h_{\text{max}} = \frac{(10)^2}{2(10)} = \frac{100}{20} = 5 \text{ m}$$

\textbf{Respuestas incorrectas:}
\begin{itemize}
    \item \textbf{B) 10 m:} Olvida dividir entre $2g$. Solo calcula $v_{0y}$ pero no aplica la fórmula completa.
    \item \textbf{C) 20 m:} Usa la velocidad total sin considerar solo la componente vertical.
    \item \textbf{D) 40 m:} Error al aplicar la fórmula. No eleva al cuadrado la velocidad vertical correctamente.
\end{itemize}
\end{respuestabox}

\begin{respuestabox}
\subsection{Pregunta 9}
\textcolor{verdecorrecto}{\textbf{Respuesta correcta: C) La aceleración gravitacional}}

\textbf{Justificación:}\\
En una gráfica velocidad-tiempo, la pendiente representa la aceleración: $\frac{dv}{dt} = a$. Para caída libre, esta pendiente es la aceleración gravitacional $g$.

\textbf{Respuestas incorrectas:}
\begin{itemize}
    \item \textbf{A) La posición del objeto:} La posición se obtiene del área bajo la curva, no de la pendiente.
    \item \textbf{B) La distancia recorrida:} La distancia es el área bajo la curva velocidad-tiempo, no la pendiente.
    \item \textbf{D) La velocidad inicial:} La velocidad inicial es el valor donde la gráfica intercepta el eje vertical en $t=0$, no la pendiente.
\end{itemize}
\end{respuestabox}

\newpage

\begin{respuestabox}
\subsection{Pregunta 10}
\textcolor{verdecorrecto}{\textbf{Respuesta correcta: B) 2 s}}

\textbf{Justificación:}\\
El tiempo de caída depende solo de la altura y la gravedad:
$$h = \frac{1}{2}gt^2$$
$$20 = \frac{1}{2}(10)t^2$$
$$20 = 5t^2$$
$$t^2 = 4$$
$$t = 2 \text{ s}$$

\textbf{Respuestas incorrectas:}
\begin{itemize}
    \item \textbf{A) 1 s:} Error al despejar. Si $t = 1$, entonces $h = \frac{1}{2}(10)(1)^2 = 5$ m, no 20 m.
    \item \textbf{C) 4 s:} Olvida extraer la raíz cuadrada. $t^2 = 4$, entonces $t = 2$, no 4.
    \item \textbf{D) 10 s:} Confunde la gravedad con el tiempo. La gravedad es un dato, no el resultado.
\end{itemize}
\end{respuestabox}

